\cotefile{157-5}
\scanpage{165}
\mspage{\unsure{81|87}}
% archive sheet number: 164
% margin note written slanted
\margin{\deleted{\unsure{Effacer|Effacé}}}

% lines 1-9
Une telle catégorie s'appellera \uline{grande devant $\pi$}.

(Il est immédiat que dans le cas où $J$ est elle-même ordonnée, cette terminologie est \unsure{cohérente} avec celle déjà adoptée pour les ens. ordonnés, i.e. que s'il existe une \unsure{application} croissante cofinale $\varphi:I\longrightarrow J$, avec $I$ grand devant $\pi$, alors $J$ est lui-$=$ grand devant $\pi$.) Une telle catégorie est automatiquement filtrante (Chap I \unsure{2.1 B b}).
% margin note written slanted
\margin{\unsure{Notons aussi} que \unsure{si on a} un foncteur cofinal $J'\longrightarrow J$, si $J'$ est grand devant $\pi$, $J$ l'est aussi (cp. \unsure{transitivité} de cofinalité).}

% lines 10-12
\uline{Proposition} 4.4.4. Soient $J$ une $\unsure{\mathcal{U}|\mathscr{U}}$-catégorie, \deleted{\unsure{filtrante}} $\pi$ un cardinal infini. Conditions équivalentes:

% lines 13-16
a) $J$ est grande devant $\pi$, i.e. il existe un ens. ordonné grand devant $\pi$ et un foncteur cofinal $\varphi:I\longrightarrow J$.

% lines 17-26
b) $J$ satisfait les deux conditions suivantes

$\mathrm{PF}_{\pi}1$: $\operatorname{Ob}J$ est grand devant $\pi$, i.e. pour toute famille $(y_{\alpha})_{\alpha\in A}$ d'objets de $J$, avec $\operatorname{card}A\leq\pi$, il existe un majorant $y$ dans $J$ (i.e. un $y$ et des flèches $y_{\alpha}\xrightarrow{v_{\alpha}}y$)

$\mathrm{PF}_{\pi}2$: Pour $y_0,y_1$ dans $J$ et \deleted{$\unsure{\forall} A$} \deleted{\unsure{pour tout} $\unsure{A}$} \deleted{deux} $(u_{\alpha})_{\alpha\in A}$ des flèches $y_0\xrightarrow{u_{\alpha}}y_1$, \inserted{\unsure{avec} $\operatorname{card}A\leq\pi$}, existe une flèche $y_1\xrightarrow{u'}y_2$ qui les égalise, i.e. telle que $u'u_{\alpha}=u'u_{\beta}$ $\forall\alpha,\beta\in A$.

(Il suffit de l'exiger pour $\operatorname{card}A\geq2$)

% margin note written slanted
\margin{Autres conditions. e) Pour $\forall$ ensemble $A$ avec $\operatorname{card}A\leq\pi$, $\delta_J:J\longrightarrow J^A$ est cofinal. e') Si $A_0$ est un ens. donné, de cardinal $=\pi$, $J\longrightarrow J^{A_0}$ est cofinal.}
% margin note written slanted
\margin{Voir aussi \unsure{4.9.6.9.1|4.4.6.9.1} et \unsure{4.9.6.9.2|4.4.6.9.2} pages \unsure{237,238}.}

% lines 27-32
c) Pour toute \deleted{\unsure{ensemble de flèches de $J$}} \inserted{partie $A$ de $\operatorname{Fl}J$}, avec $\operatorname{card}A\leq\pi$, existe une sous-catégorie $J'$ de $J$, avec $\operatorname{card}J'\leq\pi$, telle que $A\subset\operatorname{Fl}J'$ et que $J'$ admette un objet final.

% lines 33-34; sentence continues on the following scan
d) Soit $I$ l'ens des sous-catégories \inserted{$J'$} de $J$ telles que $\operatorname{card}J'\leq\pi$ et \unsure{qu'}$J'$ a un objet final,
\scanpage{166}
\mspage{82}
% archive sheet number: 165
% margin note written slanted
\margin{Introduire $I(J)$ \unsure{avant} l'énoncé de la proposition. \uline{NB} La condition de cofinalité $I(J)\longrightarrow J$ n'est pas automatique, comme on voit \unsure{par l'exemple trivial} $J=\widehat{B}_G$, $G\ne1$.}
\margin{\uline{NB} L'hyp. de cofinalité superflue si $J$ est rigide, i.e. les \unsure{endom.} d'un objet de $J$ sont identiques.}

% lines 1-5, continuation of condition d)
ordonnons $I$ par inclusion. Alors $I$ est grand devant $\pi$, de plus \deleted{\unsure{et}} \deleted{$J$ est filtrante. (Cette dernière condition est superflue si $J$ est rigide, i.e. \unsure{les} \unsure{endom.} d'un \unsure{objet de $J$} sont l'identité).} \inserted{et le foncteur canonique $I\longrightarrow J$ ($J'\longmapsto e_{J'}=$ objet final de $J'$) cofinal.}

\uline{Démonstration}

% lines 6-9
$a\Longrightarrow b$ vérification immédiate, \deleted{\unsure{en appliquant}} \inserted{en utilisant} \deleted{\unsure{de la}} condition F1, F2 de cofinalité pour \deleted{des catégories} $I$ filtrantes

$b\Longrightarrow c$ par AQT

$c\Longrightarrow d$ par AQT

% lines 12-17
$d\Longrightarrow a$ Considérons l'ens. ordonné grand devant $\pi$ construit dans l'énoncé de d), et le \uline{foncteur}
\[
I\longrightarrow J
\]
donné par
\[
J'\longmapsto e_{J'}=\text{objet final de }J'.
\]
(Vérifier que c'est un foncteur: \unsure{immédiate|immédiat}.) Ce foncteur est surjectif, donc satisfait la condition F1 de cofinalité.

% lines 18-24
La condition F2 résulte aussi de d)
\diagram{\begin{tikzcd}[ampersand replacement=\&]
j \arrow[r,shift left=1.3ex,"u_1"] \arrow[r,shift right=1.3ex,"u_2"'] \& e_{J'}
\end{tikzcd}}
$J'$, en regardant les trois objets $J'$, $j\xrightarrow{u_1}e_{J'}$, $j\xrightarrow{u_2}e_{J'}$ de $I$, et en prenant un $J''$ qui les majore tous trois, alors $e_{J'}\longrightarrow e_{J''}$ égalise $u_1,u_2$.
% margin note written slanted
\margin{Un peu bref le cas où $e_{J'}=j$ \unsure{échappe} à cet argument, et \unsure{demande} plus.}

% line 25, wholly cancelled
\deleted{\uline{Corollaire} Soit $I\longrightarrow J$ un foncteur cofinal.}

% lines 26-29
\uline{Corollaire} 4.4.5. Soit $F\in\operatorname{Ob}\widehat C$. Pour que $F$ appartienne à $\operatorname{Ind}'_{\pi}(C)$, il f. et s. que $C_{/F}$ satisfasse à l'une des quatre conditions équivalentes de 4.4.4.
% vertical marginal emphasis stroke
\margin{$\big|$}

% lines 30-33
\uline{Cor.} 4.4.6. Soit $J$ une catégorie stable par $\unsure{\underline{\lim}|\varinjlim}$ de \deleted{\unsure{type}} cardinal $\leq\pi$. Alors $J$ est grande devant $\pi$. En effet, les conditions $\mathrm{PF}_{\pi}1$, $\mathrm{PF}_{\pi}2$ sont \unsure{visiblement} vérifiées.
\scanpage{167}
\mspage{83}
% archive sheet number: 166
\deleted{Proposition}

% lines 1-7
4.4.6. Je me donne un foncteur
\[
J\xrightarrow{\varphi}\operatorname{Ind}'_{\pi}C,\tag*{(4.4.6.1)}
\]
$J$ une petite catégorie, et en cherche une $\varinjlim$ dans $\operatorname{Ind}'_{\pi}C$, en postulant au besoin l'existence des $\unsure{\underline{\lim}|\varinjlim}$ de cardinal $\leq\pi$ dans $C$, et l'hypothèse $\operatorname{card}\operatorname{Fl}J\leq\pi$.

% lines 8-14
Le procédé le plus naturel pour ramener la question de \unsure{l'existence|l'obtention} des $\underline{\lim}$ dans $\operatorname{Ind}'_{\pi}C$, à la question similaire dans $C$, est explicité dans SGA4 I 8.8.3. On considère l'inclusion pl. fid.
\[
C\overset{\alpha}{\hookrightarrow}\operatorname{Ind}'_{\pi}C,\tag*{(4.4.6.2)}
\]
laquelle commute aux $\varinjlim$ de cardinal $\leq\pi$ (celles qui sont repr. dans $C$) cf p. 78 et réf: SGA4 I 9.8), d'où
\[
\underline{\operatorname{Hom}}(J,C)\overset{\alpha}{\hookrightarrow}\underline{\operatorname{Hom}}(J,\operatorname{Ind}'_{\pi}C)\tag*{(4.4.6.3)}
\]
% lines 15-22
également pleinement fidèle, et commute aux $\varinjlim$ de cardinal $\leq\pi$ (\unsure{représentables} dans le premier \unsure{membre}). Cette \unsure{flèche} s'insère dans le carré commutatif relatif à la donnée d'un $j\in\operatorname{Ob}J$, et du foncteur d'inclusion
\[
e\xrightarrow{\varepsilon_j}J\qquad\varepsilon_j(e)=j\quad\unsure{;}\tag*{(4.4.6.4)}
\]
\scanpage{168}
\mspage{\textcircled{84}}
% archive sheet number: 167
% top diagram, the reference number is repeated as on the sheet
\diagram{\begin{tikzcd}[ampersand replacement=\&,column sep=large]
\underline{\operatorname{Hom}}(J,C) \arrow[r,"\alpha^J"] \arrow[d,"\varepsilon_j^*"] \& \underline{\operatorname{Hom}}(J,\operatorname{Ind}'_{\pi}C) \arrow[d,no head,"\varepsilon_j^{*\prime}"] \& \ni\varphi \arrow[d,mapsto] \\
C \arrow[r,"\alpha"] \& \operatorname{Ind}'_{\pi}(C) \& \ni\varphi(j)
\end{tikzcd}}
\[\tag*{(4.4.6.2)}\vphantom{C}\]
% margin note written slanted
\margin{Les foncteurs $\varepsilon_j^*$, $\varepsilon_j^{*\prime}$ sont les foncteurs évaluation en $j$.}

% lines 4-6
lequel induit un foncteur
\[
\underbrace{\underline{\operatorname{Hom}}(J,C)_{/\varphi}}_{\text{``déf'' }I}
\xrightarrow{\deleted{\unsure{\alpha}}\unsure{\varepsilon_j^{\varphi}|\xi_j^{\varphi}}}
C_{/\varphi(j)}.\tag*{(4.4.6.3)}
\]

% lines 7-19
Si $C$ est stable par $\varinjlim$ \deleted{finies, \unsure{et est}} \inserted{\unsure{relatifs|relatives} à} \unsure{un|une} $\unsure{\mathcal K|K|\mathscr K|\kappa}$ de types de \deleted{types $K$} \inserted{cardinal $\leq\pi$} (p.ex. par $\varinjlim$ finies), il en est de $=$ de $\underline{\operatorname{Hom}}(J,C)$, et aussi de $\underline{\operatorname{Hom}}(J,C)_{/\varphi}$ et de $\unsure{E|C}_{/\varphi(j)}$, en vertu du fait que $\alpha^J$ et $\alpha$ commutent aux $\varinjlim$ des types envisagés, de plus (4.4.6.3) commute aux \deleted{\unsure{Limites}} \inserted{\unsure{limites}} en question. Supposons donc $C$ \deleted{\unsure{est}} \inserted{stable} au moins par $\varinjlim$ finies, donc \deleted{\unsure{aussi}} les catégories source et but de (4.4.6.3) sont filtrantes. (Le fait que le foncteur commute aux $\varinjlim$ finies ne semble pas \unsure{utilisé|utile} dans ce contexte présent.) \deleted{\unsure{Si ce foncteur}} \inserted{\unsure{Cela peut servir|Cela après service}} Pour \unsure{voir|noter} le critère commode de cofinalité SGA4 I 8.1.3. b) par \deleted{$\unsure{\varphi}$}F1, \deleted{$\unsure{\varphi}$}F2.
\scanpage{169}
\mspage{85}
% archive sheet number: 168
% lines 1-3
Rappelons que la cofinalité des foncteurs (4.4.6.3) \inserted{quel que soit le $j\in\operatorname{Ob}J$} implique le fait que $\varphi$ est une limite inductive du foncteur \inserted{composé}

% cancelled attempts above the live diagram
\deleted{\diagram{\begin{tikzcd}[ampersand replacement=\&]
I \arrow[r,"\unsure{\psi}"] \arrow[d,equal] \& \underline{\operatorname{Hom}}(J,C) \arrow[r,"\alpha^J"] \& \unsure{\underline{\operatorname{Hom}}(J,\operatorname{Ind}'_{\pi}(C))} \\
\underline{\operatorname{Hom}}(J,C)_{/\varphi} \& \underline{\operatorname{Hom}}(J,C) \arrow[r,"\alpha^J"] \& \unsure{\underline{\operatorname{Hom}}(J,\operatorname{Ind}'_{\pi}(C))}
\end{tikzcd}}}

% main diagram
\diagram{\begin{tikzcd}[ampersand replacement=\&,column sep=large]
I \arrow[r,"\unsure{\psi^J|\psi_J}"] \arrow[d,equal] \& \underline{\operatorname{Hom}}(J,C) \arrow[r,"\alpha^J"] \arrow[dd,dotted,no head,"\varepsilon_j^*"] \& \underline{\operatorname{Hom}}(J,\operatorname{Ind}'_{\pi}(C)) \arrow[dd,dotted,"\varepsilon_j^{*\prime}"] \\
\underline{\operatorname{Hom}}(J,C)_{/\varphi} \arrow[d,"\unsure{\varepsilon_j^{\varphi}|\xi_j^{\varphi}}"'] \arrow[dr,phantom,"\deleted{\unsure{\varphi_j}}"] \& \& \\
C_{/\varphi(j)} \arrow[r,dashed,"\psi_j"] \& C \arrow[r,dashed,"\alpha"] \& \operatorname{Ind}'_{\pi}(C)
\end{tikzcd}}
\[\tag*{(4.4.6.4)}\vphantom{I}\]

\inserted{qui est l'ens. ordonné filtrant $I$ \deleted{\unsure{(plus)}} \unsure{précisément}}

% large cancelled paragraph, lines 8-13
\deleted{$\varphi$ sera $=$ la limite inductive dans la catégorie plus grande $\underline{\operatorname{Hom}}(J,\widehat C\deleted{\unsure{(C)|C}})$, admettant $\underline{\operatorname{Hom}}(J,\operatorname{Ind}_{\pi}C)$ comme sous-catégorie pleine; car \unsure{cette dernière assertion signifie}}

% lines 14-22
\unsure{plus précisément} \unsure{est que|que} $\varphi$ est la $\varinjlim$ ``argument par argument $j$'', ce qui signifie que pour tout $j\in\operatorname{Ob}J$, $\varphi(j)=\unsure{\varepsilon_j^*|\varepsilon_j^{*\prime}}(\varphi)$ est limite inductive des $\varepsilon_j^{*\prime}(\alpha^J\psi)=(\alpha\unsure{\psi_j|\varepsilon_j^*})\varepsilon_j^{\varphi}$.

\deleted{$\unsure{\varphi}$} (Or on sait \inserted{en effet} que $\varphi(j)$ est la limite de $\alpha\psi$, puisque $\varphi(j)\in\operatorname{Ob}\widehat C$ est ind-représentable, donc aussi celle des $(\alpha\psi_j)\varepsilon_j^{\varphi}$ si $\varepsilon_j^{\varphi}$ est cofinal.)

% final lines continue on next scan
Ainsi, admettant la cofinalité \inserted{de (4.4.6.3)}, \deleted{\unsure{dans}} et introduisant $\psi_J$ de (4.4.6.3)
\scanpage{170}
\mspage{86}
% archive sheet number: 169
\[
\psi_J:I\longrightarrow\underline{\operatorname{Hom}}(J,C)
\]
% lines 2-5
comme un foncteur
\[
I\times J\xrightarrow{\psi}C,
\]
on trouve la ``représentation indicielle'' de $\varphi$
\[
\varphi(j)=\underbrace{\varinjlim_{i\in I}\underbrace{\psi(i,j)}_{\deleted{\unsure{\in\operatorname{Ob}C}}}}_{\substack{\text{limite inductive}\\\text{prise dans }\operatorname{Ind}'_{\pi}(C)\\\text{d'objets de }C\subset\operatorname{Ind}'_{\pi}(C)}},
\]

% lines 6-9
d'où
\[
\begin{aligned}
\underbrace{\varinjlim_{j\in J}{}_{\operatorname{Ind}'_{\pi}C}\varphi(j)}_{\substack{\text{sous réserve}\\\text{d'existence}}}
&=\varinjlim_{j\in J}\varinjlim_{i\in I}\psi(i,j)\\
&=\varinjlim_{i\in I}\varinjlim_{j\in J}\psi(i,j),
\end{aligned}
\]
% lines 10-21
toutes les limites inductives étant prises dans $\operatorname{Ind}'_{\pi}(C)$. Mais \deleted{si \unsure{pour}} \inserted{les} $\varinjlim$ de type $J$ sont représentables dans $C$, et $\operatorname{card}J\leq\pi$, on sait que \inserted{la} $\varinjlim_{j\in J}\psi(i,j)$, \uline{dans $C$} est aussi leur $\varinjlim$ dans $\operatorname{Ind}'_{\pi}C$, et on trouve donc que $\varinjlim_{j\in J}\varphi(j)$ existe dans $\operatorname{Ind}'_{\pi}(C)$ sss $\varinjlim_{i\in I}\Psi(i)$ existe dans $\operatorname{Ind}'_{\pi}C$, où pour $i\in I$, on pose
\[
\Psi(i)=\varinjlim_{j\in J}\psi(i,j)\quad\text{\uline{dans $C$}}.
\]
\scanpage{171}
\mspage{87}
% archive sheet number: 170
% line 1
Résumons ce que nous avons obtenu:

% lines 2-12
\uline{Lemme} 4.4.7. Soit $C$ une \inserted{petite} catégorie, $\pi$ un cardinal infini, \deleted{\unsure{supposons}} $J$ une \deleted{petite} \deleted{\unsure{par}} \deleted{\unsure{type}} catégorie \deleted{de cardinal} telle que $\operatorname{card}\operatorname{Fl}J\leq\pi$, $\varphi:J\longrightarrow\operatorname{Ind}'_{\pi}(C)$ un foncteur. Supposons \inserted{les $\varinjlim$ de type $J$ existent dans $C$,} que les \deleted{\unsure{cat.}} foncteurs $\unsure{\varepsilon_j^{\varphi}|\xi_j^{\varphi}}$ ($j\in\operatorname{Ob}J$) de (4.4.6.3) soient \uline{cofinaux} (ce qui s'explicite de façon toute \unsure{simple} si source et but sont filtrants, chose acquise si $C$ est stable par $\varinjlim$ finies). Alors $\varphi$ a une ``représentation indicielle'' canonique \deleted{\unsure{par un $I$}} par
\[
\psi:I\times J\longrightarrow C\tag*{(4.4.7.1)}
\]
avec
\[
I=\underline{\operatorname{Hom}}(J,C)_{/\varphi},\tag*{(4.4.7.1)}
\]
% lines 15-17
i.e. on a l'isom. fonctoriel en $j\in\operatorname{Ob}J$
\[
\varphi\deleted{\unsure{(j)}}\simeq\varinjlim_I\psi(i,j)\tag*{(4.4.7.3)}
\]
\margin{lim. inductive dans $\operatorname{Ind}'_{\pi}(C)$ d'objets de $C$.}

% lines 18-19
Posons pour tout $i\in\operatorname{Ob}I$
\[
\Psi(i)=\varinjlim_J\psi(i,j)\tag*{(4.4.7.4)}
\]
\margin{limite inductive prise dans $C$, mais valable aussi dans $\operatorname{Ind}'_{\pi}C$}

% final lines, continuation next scan
Alors \deleted{$\varphi:J\longrightarrow\operatorname{Ind}'_{\pi}C$ a une} $\varinjlim_{j\in J}\varphi(j)$ (dans $\operatorname{Ind}'_{\pi}C$) existe sss
\scanpage{172}
\mspage{88}
% archive sheet number: 171
% lines 1-2, conclusion from previous page
la $\varinjlim_{i\in I}\Psi(i)$ existe (dans $\operatorname{Ind}'_{\pi}C$ également), et les deux sont égales.

% lines 3-11
\uline{Corollaire} \inserted{4.4.8.} Supposons que $C$ soit stable par $\varinjlim$ de cardinal $\leq\pi$ (pas seulement par celles de type $J$), \deleted{Alors \unsure{on a que}} $\underline{\operatorname{Hom}}(J,C)$ et \inserted{$I=$} $\underline{\operatorname{Hom}}(J,C)_{/\varphi}$ ont la $=$ propriété, $I$ est grand devant $\pi$, et les deux $\varinjlim_{j\in J}\varphi(j)=\varinjlim_{\unsure{j|i}\in I}\Psi(i)$ existe dans $\operatorname{Ind}'_{\pi}(C)$ (sous réserve de cofinalité des foncteurs (4.4.6.3)).

% line 12
Il suffit en effet de vérifier ceci:

% lines 13-23
\uline{Lemme} 4.4.8.1. Soit $I$ une catégorie grande devant $\pi$, et $\Psi:I\longrightarrow C$ un foncteur. Alors $\varinjlim_i\Psi(i)$ existe dans $\operatorname{Ind}'_{\pi}(C)$, et \deleted{\unsure{c'est un|est égale dans le}} \deleted{\unsure{ind-objet}} la limite inductive dans $\widehat C$ \deleted{$F$} ou encore ``c'est'' le ind-objet \inserted{grand devant $\pi$} représenté par $\Psi$. [En d'autres termes, $\operatorname{Ind}'_{\pi}(C)\subset\widehat C$ est une sous-catégorie strictement pleine \uline{stable} par $\varinjlim$ grandes devant $\pi$.]
% margin note written slanted
\margin{C'est donc cette stabilité \unsure{qui est un énoncé|est un énoncé} \unsure{nettement plus} fort! \unsure{Je} \unsure{parle prop.|par la prop.}}
\scanpage{173}
\mspage{89}
% archive sheet number: 172
% lines 1-8
Ce lemme est une tautologie: si $F$ désigne la limite inductive dans $\widehat C$, il est clair \inserted{par définition} (en prenant un foncteur $I_0\mathrel{\rule{1.5em}{0.4pt}}I$ cofinal, avec $I_0$ \uline{ordonné} grand devant $\pi$, de sorte que la $\varinjlim$ peut être prise suivant $I_0$) que $F$ est dans $\operatorname{Ind}'_{\pi}(C)$, et donc c'est à fortiori une $\varinjlim$ dans $\operatorname{Ind}'_{\pi}C$.

% lines 9-15
Pour vérifier l'existence des $\varinjlim$ de cardinal $\leq\pi$ dans $\operatorname{Ind}_{\unsure{\pi|\beta}}C$, dans le cas où $C$ lui-$=$ est stable par \unsure{petites|telles} limites, il suffit donc, grâce au corollaire 4.4.8, de prouver la cofinalité des foncteurs (4.4.6.3). Cela va résulter du \deleted{\unsure{stable par $\varinjlim$ de cardinal $\leq\pi$}}

% lines 16-26
\uline{Lemme} 4.4.\unsure{9|8}. Soient \deleted{\unsure{$\pi$}}, \inserted{$\pi$ un cardinal infini,} $C$ une petite catégorie, $J$ une petite catégorie. \deleted{Supposons $C$ stable par sommes de cardinal $\leq\pi$, et} \deleted{telle} \inserted{que} pour tous $j,j'\in\operatorname{Ob}J$, on ait $\operatorname{card}\operatorname{Hom}(j,j')\leq\pi$. \deleted{Soit alors} Dans ces conditions, \deleted{les} \inserted{pour tout $j\in\operatorname{Ob}J$,} \unsure{le} foncteur (4.4.6.3) \deleted{\unsure{est un}} \uline{admet un adjoint à gauche}, donc à fortiori il est cofinal (car les catégories colocalisées en dessous des $C_{/\varphi(j)}$ ont un objet initial, à fortiori elles sont $0$-connexes).
% margin note written slanted, crossed out
\margin{\deleted{\unsure{Adjonctions}}}
\scanpage{174}
\mspage{90}
% archive sheet number: 173
\deleted{Ce lemme est un}

% lines 1-4
Démonstration de 4.4.8. Revenons au diagramme (4.4.6.2), dont on déduit (4.4.6.3), nous allons le réécrire de façon simplifiée
\diagram{\begin{tikzcd}[ampersand replacement=\&,column sep=large]
B \arrow[r,"\beta"] \arrow[d,no head,"\varepsilon"] \& B' \arrow[d,no head,"\varepsilon'"] \& \ni y'_0 \arrow[d,mapsto] \\
C \arrow[u,bend left=45,dashed,"\lambda"] \arrow[r,no head,"\gamma"] \& C' \arrow[u,bend left=45,dashed,"\lambda'"] \& \ni z'_0
\end{tikzcd}}
\[\tag*{(4.4.9.1)}\vphantom{B}\]

% lines 5-9
Je dis que $\varepsilon,\varepsilon'$ ont \uline{des adjoints à gauche} $\lambda$ \deleted{$\unsure{\lambda'}$} resp. $\lambda'$, et que l'hom canonique de ``changement de base''
\[
\lambda'\deleted{\unsure{\beta}}\gamma\longrightarrow\beta\lambda\tag*{(4.4.9.2)}
\]
est un \uuline{iso.}

% lines 10-20
En effet, l'existence des adjoints de $\varepsilon,\varepsilon'$, à savoir des foncteurs $\varepsilon_{j!}$ et $\varepsilon'_{j!}$ relatifs à l'inclusion
\[
\varepsilon_j:e\longrightarrow J
\]
et aux catégories de foncteurs $\underline{\operatorname{Hom}}(J,C)$ resp. $\underline{\operatorname{Hom}}(J,\operatorname{Ind}'_{\pi}(C))$, \unsure{résultent|résulte} de \deleted{celle} l'existence des sommes de type $\operatorname{Hom}(j,j')$ ($j'\in\operatorname{Ob}J$) dans $C$ \inserted{(déjà acquise)} \uline{et dans} $\operatorname{Ind}'_{\pi}C$; et du fait que le foncteur $C\longrightarrow\operatorname{Ind}'_{\pi}C$ commute
\scanpage{175}
\mspage{91}
% archive sheet number: 174
% lines 1-7, continuation
auxdites sommes \deleted{\unsure{implique} \unsure{(|$C$}} (chose également \inserted{déjà} acquise, puisque $\operatorname{card}(\operatorname{Hom}(j,j'))\leq\pi$) implique donc que (4.4.9.2) est un iso. \deleted{\unsure{Il}} Moyennant la vérification en suspens (l'existence des sommes de cardinal $\leq\pi$ dans $\operatorname{Ind}'_{\pi}(C)$), le lemme 4.4.9. résulte du résultat plus général

% lines 8-16
\deleted{\uline{Corollaire}} \inserted{Lemme} 4.4.9.0. Considérons un \deleted{diagramme} \inserted{carré commutatif} (4.4.9.1) (traits pleins) de catégories, où $\varepsilon$ admet un adjoint à gauche $\lambda$, $\varepsilon'$ un adjoint à gauche $\lambda'$, et la flèche de chgt. de base (4.4.9.2) est iso. Alors pour tout $y'_0$ dans $B'_0$, posant $\varepsilon'(y'_0)=z'_0$, le foncteur induit par $(\varepsilon,\varepsilon')$:
\[
\overline B=B_{/y'_0}\xrightarrow{\overline\varepsilon}\overline C=C_{/z'_0}
\]
admet un adjoint à gauche $\overline\lambda$, donné par

% lower composite with triangle
\diagram{\begin{tikzcd}[ampersand replacement=\&,column sep=large]
{\overline\lambda\bigl(\gamma(\deleted{\unsure{y}}z)\xrightarrow{u}z'_0\bigr)=\bigl(\lambda'\gamma(z)\simeq\beta(\underline{\lambda(z)})}
\arrow[r,"\lambda'(u)"] \arrow[ddr,"\substack{\overline\lambda(u)=\text{composé}\\\mathrm{adj}\circ\lambda'(u)}"']
\& \lambda'(z'_0) \arrow[d,equal] \\
\& \lambda'\varepsilon'(y'_0) \arrow[d,"\mathrm{adj}."] \\
\& {y'_0\bigr)}
\end{tikzcd}}
\scanpage{176}
\mspage{92}
% archive sheet number: 175
% lines 1-3
La démonstration se fait par un AQT fastidieux (j'avoue ne pas avoir poussé jusqu'au dernier détail \ldots).

% lines 4-5
Pour terminer la démonstration de 4.4.\unsure{9|8}, il reste à prouver un dernier (?)

% lines 6-11, replacements above the cancelled hypotheses
\uline{Lemme 4.4.\unsure{11|10}}. Soit $C$ une \inserted{petite} catégorie, stable par \deleted{$\unsure{\varinjlim}$ \unsure{de type} $\unsure{J}$ de cardinal $\leq\pi$} \inserted{sommes de cardinal $\leq\pi$}. Alors $\operatorname{Ind}'_{\pi}C$ est stable \deleted{par les \unsure{mêmes} \unsure{sommes}} \inserted{\unsure{sommes} de cardinal $\leq\pi$} (et, pour mémoire, le foncteur $\alpha:C\mathrel{\rule{1.5em}{0.4pt}}\operatorname{Ind}'_{\pi}C$ y commute).

% lines 12-21
Appliquant \deleted{le lemme} \inserted{le cor.} 4.4.\unsure{8|9} ou le lemme 4.4.7, on est ramené à prouver la cofinalité des (4.4.6.3), quand $J$ est discret, $+$ le fait que $\underline{\operatorname{Hom}}(J,C)_{/\varphi}$ est grande devant $\pi$, dans ce cas. Mais la donnée de $\varphi$ revient à celle d'une famille $(\varphi_j)_{j\in J}$ d'objets de $\operatorname{Ind}'_{\pi}(C)$, et les foncteurs (4.4.6.3) \unsure{s'explicitent|s'explicite} comme les projections
\[
\prod_{j\in J}(C_{/\varphi_j})\unsure{\longleftrightarrow|\longrightarrow} C_{/\varphi_j}
\]
du produit total sur un des facteurs. Donc on est ramené à ceci:
\scanpage{177}
\mspage{93}
% archive sheet number: 176
% lines 1-6
\uline{Lemme} 4.4.12. Soit \deleted{\unsure{les}} $(C_j)_{j\in J}$ une famille de \inserted{petites} catégories grandes devant $\pi$, indexée par un \inserted{petit} ensemble $J$ \deleted{tel que $\operatorname{card}J\leq\pi$}. Alors $\prod C_j$ est grande devant $\pi$, et les projections (pour $j_0\in J$) $\prod C_j\longrightarrow C_{j_0}$ sont cofinales.

% lines 7-16
Prouvons d'abord la cofinalité, qui est la plus importante pour nous. C'est \unsure{évident} si $J=\{j_0\}$, sinon on a $J=J'\cup\{j_0\}$, et $\prod_{j\in J}C_j\simeq C'\times C_{j_0}$, où $C'=\prod_{j\in J'}C_j$, donc il faut voir que $C'\times C_{j_0}\longrightarrow C_{j_0}$ est cofinal, i.e. que les catégories ${}_{\xi}\backslash(C'\times C_{j_0})$ (pour $\xi\in\operatorname{Ob}C_{j_0}$) sont $0$-connexes, or \deleted{elles sont} elles se \unsure{ramènent|réduisent}: $C'\times{}_{\xi}\backslash C_{j_0}$, et ${}_{\xi}\backslash C_{j_0}$ est $0$-connexe, il reste à voir que $C'$ est $0$-connexe.

% lines 17-24
Or \unsure{faire|fera} attention à cet égard qu'un produit \uline{infini} de catégories $0$-connexes n'est pas forcément $0$-connexe, donc \deleted{\unsure{il}} la \deleted{elle n'est pas} \inserted{cofinalité} n'est pas totalement tautologique! Mais si on prouve que la catégorie $\prod_{j\in C}C_j$ est
\scanpage{178}
\mspage{94}
% archive sheet number: 177
% lines 1-3, continuation
\unsure{grde|grande} devant $\pi$ donc filtrante, cela implique qu'elle est $0$-connexe, et on peut appliquer le $=$ résultat à $\inserted{C'=}\prod_{j\in J'}C_j$, et on gagne.

% lines 4-12
Donc on est ramené à prouver que \deleted{\unsure{le}} $\prod C_j$ est grand devant $\pi$. Ceci \deleted{peut} se voit soit sur le critère $\mathrm{PF}_{\pi}1,\mathrm{PF}_{\pi}2$ de 4.4.4 b, soit \deleted{sur la définition} en se ramenant au cas \deleted{des} où les $C_j$ sont \unsure{ordonnées|ordonnés}, mais cela a l'air plus long. Notons que l'on prouve de $=$ qu'un \inserted{(petit)} produit de catégories filtrantes est encore \inserted{une cat.} filtrante.

% lines 13-20
\uline{Remarques} 1. Cet argument ne montre nullement ($=$ si $C$ est stable par $\unsure{\varinjlim|\underline{\lim}}$ de cardinal $\leq\pi$) que pour tout $\varphi\in\underline{\operatorname{Hom}}(J,\operatorname{Ind}(\deleted{\unsure{J|j}}C))$, le foncteur induit:
\[
\underline{\operatorname{Hom}}(\unsure{I|J},C)_{/\varphi}\longrightarrow C_{/\varphi(j)}
\tag*{(4.4.6.3)}
\]
est cofinal, et je doute qu'il le soit toujours. La difficulté provient du fait que $C\longrightarrow\operatorname{Ind}C$ ne commute pas aux sommes de cardinal $\leq\pi$, seulement aux sommes finies.
\scanpage{179}
\mspage{95}
% archive sheet number: 178
% line 1, continuation
donc le lemme 4.4.10 ne peut être invoqué.

% lines 2-12
2. Nous venons de prouver, via 4.4.8 et 4.4.9 (dont la démonstration est achevée avec 4.4.12), que \deleted{\unsure{le}} si $C$ est une petite catégorie stable par $\unsure{\varinjlim|\underline{\lim}}$ de cardinal $\leq\pi$, alors $\operatorname{Ind}_{\pi}C$ est également stable par ces limites. \deleted{Donc} \inserted{De plus} \deleted{pour} les \deleted{limites de type $J$} systèmes inductifs de type $J$ dans $\operatorname{Ind}_{\pi}C$ \deleted{\unsure{rep|resp.|se}} admettent une représentation indicielle
\[
\deleted{\unsure{\varphi}}\quad\psi:I\times J\mathrel{\rule{1.5em}{0.4pt}}C,
\]
avec $I$ grand devant $\pi$. Et on peut toujours supposer, quitte à prendre $I'\longrightarrow I$ cofinal avec $I'$ ens. ordonné grand devant $\pi$, que $I$ est \inserted{de plus} un ens. ordonné.

% lines 18-20
C'est là le \unsure{noyau} \unsure{dur|du}, pour la démonstration du

% lines 21-25
\uline{Théorème 4.4.13.} Soit $C$ une catégorie équivalente à une petite catégorie, et stable par $\varinjlim$ de cardinal $\leq\pi$, où $\pi$ est un cardinal infini donné. Alors: \deleted{$\operatorname{Ind}_{\pi}$ \unsure{jouit|pour} \illegible{1}}
\scanpage{180}
\mspage{96}
% archive sheet number: 179
% lines 1-3
a) (pour mémoire) le foncteur canonique
\[
C\longrightarrow\operatorname{Ind}_{\pi}(C)\unsure{\xrightarrow{\approx}|\approxeq|\simeq|\cong}\operatorname{Ind}'_{\pi}(C)
\]
commute aux $\unsure{\underline{\lim}|\varinjlim}$ de cardinal $\leq\pi$, et

% lines 4-11
b) (pour mémoire) les systèmes inductifs de cardinal $\leq\pi$ dans $\operatorname{Ind}_{\pi}(C)$ (i.e. les foncteurs $\varphi:J\mathrel{\rule{1.5em}{0.4pt}}\operatorname{Ind}_{\pi}(C)$, avec $\operatorname{card}\operatorname{Fl}J\leq\pi$) admettent une représentation indicielle
\[
\psi:I\times J\longrightarrow C,
\]
avec $I$ ens. ordonné grand devant $\pi$, et admettent une limite inductive dans $\operatorname{Ind}_{\pi}C$ qui se calcule par
\[
\underbrace{\varinjlim_{j\in J}\varphi(j)}_{\text{dans }\operatorname{Ind}'_{\pi}(C)}\simeq\text{``}\varinjlim_{i\in I}\text{''}\Psi(i)
\]
où
\[
\Psi(i)=\varinjlim_{j\in J}{}_{C}\psi(i,j).
\]

% lines 17-18
c) $\operatorname{Ind}_{\pi}C$ est stable par petites limites inductives.

% margin note written sideways
\margin{\unsure{rajouter} ici, en d), le cor. 4.4.\unsure{16|15} \unsure{pge|page} 101}

% bottom paragraph
Il reste à établir ce dernier point. Comme on a déjà les $\underline{\lim}$ finies par b), il reste à prouver l'existence des \inserted{petites} $\underline{\lim}$ filtrantes, \deleted{\unsure{des} $\underline{\lim}$ $\unsure{I}$ \unsure{quelconques}} \inserted{\unsure{ou} \unsure{en} \illegible{1}} \deleted{filtrantes} des sommes infinies, ce qui est plus facile en fait. Mais
